% !TEX program = pdflatex
\documentclass[11pt,a4paper]{article}

% =========================================================================
% CRONOGRAMA DEL PROYECTO
% Proyecto: CasaScore (AlquilaPuno)
% Gestión de Proyectos TICS (EST315) -- Semestre 2026-II -- Team Sobrevivientes
% =========================================================================

% --- PAQUETES DE CONFIGURACIÓN ---
\usepackage[utf8]{inputenc}
\usepackage[T1]{fontenc}
\usepackage[spanish,es-tabla]{babel}
% Compatibilidad con babel-spanish reciente (flechas de TikZ y signo de porcentaje)
\AtBeginDocument{\shorthandoff{<>}\spanishplainpercent}
\usepackage{geometry}
\geometry{
    top=3.6cm,
    bottom=2.3cm,
    left=2.2cm,
    right=2.2cm,
    % PARTE MODIFICABLE: headheight debe ser mayor que la altura del logo más alto (2.56cm = 73pt).
    headheight=80pt,
    headsep=15pt,
    footskip=20pt
}

\usepackage{graphicx}
\usepackage{fancyhdr}
\usepackage[table]{xcolor}
\usepackage{titlesec}
\usepackage{booktabs}
\usepackage{tabularx}
\usepackage{array}
\usepackage{amsmath}
\usepackage{enumitem}
\usepackage{microtype}
\usepackage{amssymb}
% --- RECURSOS VISUALES: GRÁFICOS Y CÓDIGO DE COLORES ---
\usepackage{tikz}
\usetikzlibrary{positioning,calc}
\definecolor{GrisFila}{RGB}{246, 248, 251}
\definecolor{CuadCerca}{RGB}{150, 30, 30}
\definecolor{CuadSatisf}{RGB}{180, 130, 20}
\definecolor{CuadInfo}{RGB}{10, 45, 90}
\definecolor{CuadMonit}{RGB}{110, 115, 125}
\newcommand{\etiqueta}[2]{{\setlength{\fboxsep}{2pt}\colorbox{#1}{\textcolor{white}{\scriptsize\bfseries #2}}}}
\newcommand{\cerca}{\etiqueta{CuadCerca}{Gestionar de cerca}}
\newcommand{\satisf}{\etiqueta{CuadSatisf}{Mantener satisfecho}}
\newcommand{\informado}{\etiqueta{CuadInfo}{Mantener informado}}
\newcommand{\monit}{\etiqueta{CuadMonit}{Monitorear}}
\usepackage{needspace}
\usepackage{longtable}
\definecolor{Rama1}{RGB}{10, 45, 90}
\definecolor{Rama2}{RGB}{40, 95, 150}
\definecolor{Rama3}{RGB}{180, 130, 20}
\definecolor{Rama4}{RGB}{30, 110, 60}
\definecolor{Rama5}{RGB}{110, 60, 140}
\definecolor{RaciR}{RGB}{150, 30, 30}
\definecolor{RaciA}{RGB}{10, 45, 90}
\definecolor{RaciC}{RGB}{180, 130, 20}
\definecolor{RaciI}{RGB}{140, 145, 155}
\definecolor{NivAlto}{RGB}{30, 110, 60}
\definecolor{NivMedio}{RGB}{180, 130, 20}
\definecolor{NivBajo}{RGB}{150, 30, 30}
\usetikzlibrary{patterns}

\usepackage{hyperref}

% Logos: funciona a 2 o 3 niveles de profundidad dentro de docs/latex/
\graphicspath{{../../assets/img/}{../../../assets/img/}}

% --- PALETA INSTITUCIONAL FIEI / UNA PUNO ---
\definecolor{AzulUNA}{RGB}{10, 45, 90}
\definecolor{DoradoFIEI}{RGB}{180, 130, 20}
\definecolor{GrisOscuro}{RGB}{50, 50, 50}
\definecolor{AzulClaro}{RGB}{228, 236, 246}

% --- CONFIGURACIÓN DE LISTAS ---
\setlist{topsep=2pt, itemsep=2pt, parsep=1pt}

% --- COLUMNAS DE TABLA ---
\newcolumntype{L}{>{\raggedright\arraybackslash}X}
\newcolumntype{P}[1]{>{\raggedright\arraybackslash}p{#1}}
\newcolumntype{C}[1]{>{\centering\arraybackslash}p{#1}}
\newcolumntype{Q}[1]{>{\raggedleft\arraybackslash}p{#1}}
\newcolumntype{R}{>{\raggedleft\arraybackslash}X}
\renewcommand{\arraystretch}{1.25}

% --- FORMATO DE ENLACES ---
\hypersetup{
    colorlinks=true,
    linkcolor=AzulUNA,
    urlcolor=AzulUNA,
    citecolor=AzulUNA,
    pdftitle={CRONOGRAMA DEL PROYECTO -- CasaScore (AlquilaPuno)},
    pdfauthor={Team Sobrevivientes}
}

% --- FORMATO DE TÍTULOS ---
\titleformat{\section}{\color{AzulUNA}\normalfont\large\bfseries}{\thesection.}{0.7em}{}[\color{AzulUNA}\titlerule]
\titleformat{\subsection}{\color{AzulUNA}\normalfont\normalsize\bfseries}{\thesubsection.}{0.5em}{}
\titlespacing*{\section}{0pt}{9pt}{4pt}
\titlespacing*{\subsection}{0pt}{6pt}{2pt}

% --- ENCABEZADO Y PIE DE PÁGINA (fancyhdr) ---
\pagestyle{fancy}
\fancyhf{}

% ===== TAMAÑO DE LOS LOGOS (PARTE MODIFICABLE) =====
% logo_fiei.png  : height=1.45cm (lado izquierdo)
% logo_equipo.png: height=2.56cm (lado derecho; 25 % más grande porque la imagen tiene más borde blanco)
% Para cambiar el tamaño, modifique solo el valor height=... de cada \includegraphics.
% Si el logo más alto supera 2.6cm, aumente también headheight y top en la configuración de página.
\fancyhead[L]{%
    \includegraphics[height=1.45cm,keepaspectratio]{logo_fiei.png}%
}
\fancyhead[C]{%
    \scriptsize \textbf{\textcolor{AzulUNA}{UNIVERSIDAD NACIONAL DEL ALTIPLANO DE PUNO}}\\[2pt]
    \tiny \textcolor{GrisOscuro}{FACULTAD DE INGENIERÍA ESTADÍSTICA E INFORMÁTICA}\\[2pt]
    \scriptsize \textbf{\textcolor{AzulUNA}{Gestión de Proyectos TICS (EST315) \textbar\ Team Sobrevivientes}}\\[2pt]
    \footnotesize \textbf{\textcolor{DoradoFIEI}{CRONOGRAMA DEL PROYECTO}}%
}
\fancyhead[R]{%
    \includegraphics[height=2.56cm,keepaspectratio]{logo_equipo.png}%
}

\renewcommand{\headrulewidth}{0pt}
\renewcommand{\headrule}{}

\fancyfoot[L]{\scriptsize \textcolor{GrisOscuro}{CasaScore (AlquilaPuno)}}
\fancyfoot[C]{\small \thepage}
\fancyfoot[R]{\scriptsize \textcolor{GrisOscuro}{Puno, Perú -- Semestre 2026-II}}
\renewcommand{\footrulewidth}{0.4pt}

\begin{document}

\begin{center}
    {\Large \textbf{\textcolor{AzulUNA}{CRONOGRAMA DEL PROYECTO}}}\\[3pt]
    {\normalsize \textbf{CasaScore (AlquilaPuno)}}\\[4pt]
    \footnotesize \textbf{Fecha:} viernes 18/09/2026 (semana 4) \quad \textbar\quad \textbf{Versión:} 1.0 \quad \textbar\quad \textbf{Fase:} Planificación
\end{center}

\vspace{0.1cm}

\section{Datos Generales}
\noindent\begin{tabularx}{\textwidth}{@{}P{4.2cm} L@{}}
\toprule
\textbf{Proyecto:} & CasaScore (AlquilaPuno) \\
\textbf{Responsable:} & Ruth Karina Apaza Solis -- Analista de Planificación \\
\textbf{Proceso PM4DEV:} & Gestión del cronograma \\
\textbf{Insumos:} & EDT y diccionario de la EDT, plan del equipo y Acta de Constitución \\
\textbf{Dedicación:} & 8 horas semanales por integrante (32 horas del equipo por semana) \\
\bottomrule
\end{tabularx}

\section{Método}
\begin{itemize}[leftmargin=*]
    \item \textbf{Actividades:} cada paquete de trabajo de la EDT es una actividad del cronograma. Las duraciones están en semanas de calendario, con la dedicación parcial del equipo.
    \item \textbf{Estimación PERT:} para cada actividad se estima una duración optimista ($o = 0.8\,m$), más probable ($m$) y pesimista ($p = 1.5\,m$). La duración esperada es $t_e = (o + 4m + p)/6$ y la desviación es $\sigma = (p - o)/6$. Las actividades terminadas usan su duración real.
    \item \textbf{Ruta crítica:} se calcula con inicio y fin tempranos (pasada hacia adelante) y tardíos (pasada hacia atrás). Las actividades con holgura cero forman la ruta crítica.
    \item \textbf{Calendario:} la semana 1 del semestre no se cuenta. En las semanas 9 (exámenes parciales) y 17 (exámenes finales) solo avanzan las actividades de gestión y documentación; el desarrollo, las pruebas y el trabajo de campo se detienen.
\end{itemize}

\section{Calendario del Proyecto}
{\small
\setlength{\tabcolsep}{4pt}
\noindent\begin{tabularx}{\textwidth}{@{}C{1.0cm} L C{1.0cm} L C{1.0cm} L@{}}
\toprule
\rowcolor{AzulClaro}
\textbf{Sem.} & \textbf{Fechas} & \textbf{Sem.} & \textbf{Fechas} & \textbf{Sem.} & \textbf{Fechas} \\
\midrule
2 & 31/08 -- 04/09 & 8 & 12/10 -- 16/10 & 14 & 23/11 -- 27/11 \\
3 & 07/09 -- 11/09 & 9 & 19/10 -- 23/10 \textcolor{RaciR}{\scriptsize(Exámenes parciales)} & 15 & 30/11 -- 04/12 \\
4 & 14/09 -- 18/09 & 10 & 26/10 -- 30/10 & 16 & 07/12 -- 11/12 \\
5 & 21/09 -- 25/09 & 11 & 02/11 -- 06/11 & 17 & 14/12 -- 18/12 \textcolor{RaciR}{\scriptsize(Exámenes finales)} \\
6 & 28/09 -- 02/10 & 12 & 09/11 -- 13/11 & 18 & \textcolor{GrisOscuro}{21/12 -- 25/12 \textcolor{RaciR}{\scriptsize(Cierre del semestre (24/12))}} \\
7 & 05/10 -- 09/10 & 13 & 16/11 -- 20/11 &  &  \\
\bottomrule
\end{tabularx}
}

\Needspace{14\baselineskip}
\section{Red de Actividades}
Cada nodo muestra el código de la actividad y su duración esperada en semanas. Los nodos se ordenan de izquierda a derecha según sus predecesoras y los de borde rojo forman la ruta crítica.

\begin{center}
\begin{tikzpicture}[font=\tiny, nd/.style={draw=#1, fill=#1!10, line width=0.6pt, rounded corners=1.5pt, minimum width=1.0cm, minimum height=0.62cm, align=center, inner sep=1pt}, >=stealth]
\node[nd=Rama1] (n1_1) at (0.60,-1.05) {\textbf{1.1}\\ $t_e$ 2.00};
\node[nd=RaciR, line width=1.2pt] (n1_2) at (1.96,0.00) {\textbf{1.2}\\ $t_e$ 1.00};
\node[nd=Rama1] (n1_3) at (3.33,-3.15) {\textbf{1.3}\\ $t_e$ 1.05};
\node[nd=Rama1] (n1_4) at (4.69,-1.05) {\textbf{1.4}\\ $t_e$ 1.05};
\node[nd=Rama1] (n1_6) at (7.42,-1.05) {\textbf{1.6}\\ $t_e$ 1.89};
\node[nd=RaciR, line width=1.2pt] (n2_1) at (0.60,0.00) {\textbf{2.1}\\ $t_e$ 1.00};
\node[nd=RaciR, line width=1.2pt] (n2_2) at (3.33,0.00) {\textbf{2.2}\\ $t_e$ 1.05};
\node[nd=RaciR, line width=1.2pt] (n2_3) at (3.33,-1.05) {\textbf{2.3}\\ $t_e$ 1.05};
\node[nd=RaciR, line width=1.2pt] (n3_1) at (3.33,-2.10) {\textbf{3.1}\\ $t_e$ 2.10};
\node[nd=RaciR, line width=1.2pt] (n3_2) at (4.69,0.00) {\textbf{3.2}\\ $t_e$ 1.05};
\node[nd=RaciR, line width=1.2pt] (n3_3) at (6.05,0.00) {\textbf{3.3}\\ $t_e$ 2.10};
\node[nd=RaciR, line width=1.2pt] (n3_4) at (7.42,0.00) {\textbf{3.4}\\ $t_e$ 3.15};
\node[nd=Rama3] (n3_5) at (7.42,-2.10) {\textbf{3.5}\\ $t_e$ 2.10};
\node[nd=RaciR, line width=1.2pt] (n3_6) at (8.78,-1.05) {\textbf{3.6}\\ $t_e$ 1.05};
\node[nd=Rama4] (n4_1) at (6.05,-1.05) {\textbf{4.1}\\ $t_e$ 1.05};
\node[nd=Rama4] (n4_2) at (4.69,-2.10) {\textbf{4.2}\\ $t_e$ 1.05};
\node[nd=RaciR, line width=1.2pt] (n4_3) at (8.78,0.00) {\textbf{4.3}\\ $t_e$ 1.05};
\node[nd=RaciR, line width=1.2pt] (n4_4) at (10.15,0.00) {\textbf{4.4}\\ $t_e$ 1.05};
\node[nd=RaciR, line width=1.2pt] (n5_1) at (11.51,0.00) {\textbf{5.1}\\ $t_e$ 1.05};
\node[nd=RaciR, line width=1.2pt] (n5_2) at (12.87,0.00) {\textbf{5.2}\\ $t_e$ 1.57};
\node[nd=RaciR, line width=1.2pt] (n5_3) at (14.24,0.00) {\textbf{5.3}\\ $t_e$ 0.53};
\node[draw=AzulUNA, fill=AzulUNA, text=white, circle, inner sep=1pt, font=\tiny\bfseries] (ini) at (-0.55,0) {I};
\node[draw=AzulUNA, fill=AzulUNA, text=white, circle, inner sep=1pt, font=\tiny\bfseries] (fin) at (15.19,0) {F};
\draw[->, gray!70] (ini) to[out=0,in=180] (n1_1.west);
\draw[->, RaciR, thick] (n2_1.east) to[out=0,in=180] (n1_2.west);
\draw[->, gray!60] (n1_2.east) to[out=0,in=180] (n1_3.west);
\draw[->, gray!60] (n1_3.east) to[out=0,in=180] (n1_4.west);
\draw[->, gray!60] (n3_3.east) to[out=0,in=180] (n1_6.west);
\draw[->, RaciR, thick] (ini) to[out=0,in=180] (n2_1.west);
\draw[->, RaciR, thick] (n1_2.east) to[out=0,in=180] (n2_2.west);
\draw[->, RaciR, thick] (n1_2.east) to[out=0,in=180] (n2_3.west);
\draw[->, RaciR, thick] (n1_2.east) to[out=0,in=180] (n3_1.west);
\draw[->, RaciR, thick] (n2_2.east) to[out=0,in=180] (n3_2.west);
\draw[->, RaciR, thick] (n2_3.east) to[out=0,in=180] (n3_2.west);
\draw[->, RaciR, thick] (n3_1.east) to[out=0,in=180] (n3_3.west);
\draw[->, RaciR, thick] (n3_2.east) to[out=0,in=180] (n3_3.west);
\draw[->, RaciR, thick] (n3_3.east) to[out=0,in=180] (n3_4.west);
\draw[->, gray!60] (n3_3.east) to[out=0,in=180] (n3_5.west);
\draw[->, RaciR, thick] (n3_4.east) to[out=0,in=180] (n3_6.west);
\draw[->, gray!60] (n3_5.east) to[out=0,in=180] (n3_6.west);
\draw[->, gray!60] (n3_1.east) to[out=0,in=180] (n4_1.west);
\draw[->, gray!60] (n3_2.east) to[out=0,in=180] (n4_1.west);
\draw[->, gray!60] (n2_3.east) to[out=0,in=180] (n4_2.west);
\draw[->, RaciR, thick] (n3_4.east) to[out=0,in=180] (n4_3.west);
\draw[->, gray!60] (n3_5.east) to[out=0,in=180] (n4_3.west);
\draw[->, gray!60] (n4_2.east) to[out=0,in=180] (n4_3.west);
\draw[->, RaciR, thick] (n3_6.east) to[out=0,in=180] (n4_4.west);
\draw[->, RaciR, thick] (n4_3.east) to[out=0,in=180] (n4_4.west);
\draw[->, gray!60] (n4_1.east) to[out=0,in=180] (n4_4.west);
\draw[->, RaciR, thick] (n4_4.east) to[out=0,in=180] (n5_1.west);
\draw[->, RaciR, thick] (n5_1.east) to[out=0,in=180] (n5_2.west);
\draw[->, RaciR, thick] (n5_2.east) to[out=0,in=180] (n5_3.west);
\draw[->, gray!70] (n1_1.east) to[out=0,in=180] (fin);
\draw[->, gray!70] (n1_4.east) to[out=0,in=180] (fin);
\draw[->, gray!70] (n1_6.east) to[out=0,in=180] (fin);
\draw[->, RaciR, thick] (n5_3.east) to[out=0,in=180] (fin);
\end{tikzpicture}
\end{center}

\section{Estimación PERT}
Duraciones en semanas. Las filas resaltadas pertenecen a la ruta crítica.

\vspace{0.1cm}
{\small
\setlength{\tabcolsep}{3pt}
\begin{longtable}{@{}P{0.75cm} >{\raggedright\arraybackslash}p{4.3cm} P{1.95cm} Q{0.85cm} Q{0.85cm} Q{0.85cm} Q{0.95cm} Q{0.85cm} C{1.75cm}@{}}
\toprule
\rowcolor{AzulClaro}
\textbf{Cód.} & \textbf{Actividad} & \textbf{Pred.} & \textbf{$o$} & \textbf{$m$} & \textbf{$p$} & \textbf{$t_e$} & \textbf{$\sigma$} & \textbf{Estado} \\
\midrule
\endfirsthead
\toprule
\rowcolor{AzulClaro}
\textbf{Cód.} & \textbf{Actividad} & \textbf{Pred.} & \textbf{$o$} & \textbf{$m$} & \textbf{$p$} & \textbf{$t_e$} & \textbf{$\sigma$} & \textbf{Estado} \\
\midrule
\endhead
\bottomrule
\endfoot
1.1 & Gobernanza del equipo & Inicio & 2.00 & 2.00 & 2.00 & \textbf{2.00} & 0.00 & Terminada \\
\rowcolor{RaciR!8}\textbf{1.2} & Viabilidad y Acta de Constitución & 2.1 & 1.00 & 1.00 & 1.00 & \textbf{1.00} & 0.00 & Terminada \\
1.3 & Planificación base & 1.2 & 0.80 & 1.00 & 1.50 & \textbf{1.05} & 0.12 & En curso \\
1.4 & Planes de soporte & 1.3 & 0.80 & 1.00 & 1.50 & \textbf{1.05} & 0.12 & Por iniciar \\
1.5 & Seguimiento y control & 1.4 & \multicolumn{5}{c}{\footnotesize nivel de esfuerzo: dura todo el proyecto} & En curso \\
1.6 & Entregas de la Unidad II & 3.3 & 1.44 & 1.80 & 2.70 & \textbf{1.89} & 0.21 & Por iniciar \\
\rowcolor{RaciR!8}\textbf{2.1} & Ficha técnica y requerimientos & Inicio & 1.00 & 1.00 & 1.00 & \textbf{1.00} & 0.00 & Terminada \\
\rowcolor{RaciR!8}\textbf{2.2} & Modelo de datos y arquitectura & 1.2 & 0.80 & 1.00 & 1.50 & \textbf{1.05} & 0.12 & En curso \\
\rowcolor{RaciR!8}\textbf{2.3} & Diseño de pantallas & 1.2 & 0.80 & 1.00 & 1.50 & \textbf{1.05} & 0.12 & En curso \\
\rowcolor{RaciR!8}\textbf{3.1} & Motor CASA SCORE & 1.2 & 1.60 & 2.00 & 3.00 & \textbf{2.10} & 0.23 & En curso \\
\rowcolor{RaciR!8}\textbf{3.2} & Publicación de cuartos & 2.2, 2.3 & 0.80 & 1.00 & 1.50 & \textbf{1.05} & 0.12 & Por iniciar \\
\rowcolor{RaciR!8}\textbf{3.3} & Ranking y detalle (MVP Unidad I) & 3.1, 3.2 & 1.60 & 2.00 & 3.00 & \textbf{2.10} & 0.23 & Por iniciar \\
\rowcolor{RaciR!8}\textbf{3.4} & Nube, registro e inicio de sesión & 3.3 & 2.40 & 3.00 & 4.50 & \textbf{3.15} & 0.35 & Por iniciar \\
3.5 & Punto de referencia y nuevos criterios & 3.3 & 1.60 & 2.00 & 3.00 & \textbf{2.10} & 0.23 & Por iniciar \\
\rowcolor{RaciR!8}\textbf{3.6} & Perfil y mis cuartos & 3.4, 3.5 & 0.80 & 1.00 & 1.50 & \textbf{1.05} & 0.12 & Por iniciar \\
4.1 & Pruebas del motor & 3.1, 3.2 & 0.80 & 1.00 & 1.50 & \textbf{1.05} & 0.12 & Por iniciar \\
4.2 & Levantamiento de cuartos & 2.3 & 0.80 & 1.00 & 1.50 & \textbf{1.05} & 0.12 & Por iniciar \\
\rowcolor{RaciR!8}\textbf{4.3} & Piloto con clientes & 3.4, 3.5, 4.2 & 0.80 & 1.00 & 1.50 & \textbf{1.05} & 0.12 & Por iniciar \\
\rowcolor{RaciR!8}\textbf{4.4} & Pruebas de aceptación & 3.6, 4.3, 4.1 & 0.80 & 1.00 & 1.50 & \textbf{1.05} & 0.12 & Por iniciar \\
\rowcolor{RaciR!8}\textbf{5.1} & Manuales & 4.4 & 0.80 & 1.00 & 1.50 & \textbf{1.05} & 0.12 & Por iniciar \\
\rowcolor{RaciR!8}\textbf{5.2} & Informe final y exposición & 5.1 & 1.20 & 1.50 & 2.25 & \textbf{1.57} & 0.17 & Por iniciar \\
\rowcolor{RaciR!8}\textbf{5.3} & Lecciones aprendidas y acta de cierre & 5.2 & 0.40 & 0.50 & 0.75 & \textbf{0.53} & 0.06 & Por iniciar \\
\end{longtable}
}

\Needspace{12\baselineskip}
\section{Ruta Crítica}
{\small
\setlength{\tabcolsep}{3pt}
\begin{longtable}{@{}P{0.75cm} >{\raggedright\arraybackslash}p{5.2cm} C{1.6cm} C{1.6cm} C{1.6cm} C{1.6cm} Q{1.3cm}@{}}
\toprule
\rowcolor{AzulClaro}
\textbf{Cód.} & \textbf{Actividad} & \textbf{Inicio temprano} & \textbf{Fin temprano} & \textbf{Inicio tardío} & \textbf{Fin tardío} & \textbf{Holgura (sem.)} \\
\midrule
\endfirsthead
\toprule
\rowcolor{AzulClaro}
\textbf{Cód.} & \textbf{Actividad} & \textbf{Inicio temprano} & \textbf{Fin temprano} & \textbf{Inicio tardío} & \textbf{Fin tardío} & \textbf{Holgura (sem.)} \\
\midrule
\endhead
\bottomrule
\endfoot
1.1 & Gobernanza del equipo & 31/08 & 11/09 & 03/12 & 16/12 & 13.60 \\
\rowcolor{RaciR!8}\textbf{1.2} & Viabilidad y Acta de Constitución & 07/09 & 11/09 & 07/09 & 11/09 & \textcolor{RaciR}{\textbf{0.00}} \\
1.3 & Planificación base & 14/09 & 21/09 & 02/12 & 09/12 & 11.50 \\
1.4 & Planes de soporte & 21/09 & 28/09 & 09/12 & 16/12 & 11.50 \\
1.5 & Seguimiento y control & 28/09 & 16/12 & 28/09 & 16/12 & -- \\
1.6 & Entregas de la Unidad II & 13/10 & 26/10 & 03/12 & 16/12 & 7.51 \\
\rowcolor{RaciR!8}\textbf{2.1} & Ficha técnica y requerimientos & 31/08 & 04/09 & 31/08 & 04/09 & \textcolor{RaciR}{\textbf{0.00}} \\
\rowcolor{RaciR!8}\textbf{2.2} & Modelo de datos y arquitectura & 14/09 & 21/09 & 14/09 & 21/09 & \textcolor{RaciR}{\textbf{0.00}} \\
\rowcolor{RaciR!8}\textbf{2.3} & Diseño de pantallas & 14/09 & 21/09 & 14/09 & 21/09 & \textcolor{RaciR}{\textbf{0.00}} \\
\rowcolor{RaciR!8}\textbf{3.1} & Motor CASA SCORE & 14/09 & 28/09 & 14/09 & 28/09 & \textcolor{RaciR}{\textbf{0.00}} \\
\rowcolor{RaciR!8}\textbf{3.2} & Publicación de cuartos & 21/09 & 28/09 & 21/09 & 28/09 & \textcolor{RaciR}{\textbf{0.00}} \\
\rowcolor{RaciR!8}\textbf{3.3} & Ranking y detalle (MVP Unidad I) & 28/09 & 12/10 & 28/09 & 12/10 & \textcolor{RaciR}{\textbf{0.00}} \\
\rowcolor{RaciR!8}\textbf{3.4} & Nube, registro e inicio de sesión & 13/10 & 10/11 & 13/10 & 10/11 & \textcolor{RaciR}{\textbf{0.00}} \\
3.5 & Punto de referencia y nuevos criterios & 13/10 & 03/11 & 27/10 & 10/11 & 2.05 \\
\rowcolor{RaciR!8}\textbf{3.6} & Perfil y mis cuartos & 10/11 & 17/11 & 10/11 & 17/11 & \textcolor{RaciR}{\textbf{0.00}} \\
4.1 & Pruebas del motor & 28/09 & 05/10 & 10/11 & 17/11 & 6.25 \\
4.2 & Levantamiento de cuartos & 21/09 & 28/09 & 03/11 & 10/11 & 6.25 \\
\rowcolor{RaciR!8}\textbf{4.3} & Piloto con clientes & 10/11 & 17/11 & 10/11 & 17/11 & \textcolor{RaciR}{\textbf{0.00}} \\
\rowcolor{RaciR!8}\textbf{4.4} & Pruebas de aceptación & 18/11 & 25/11 & 18/11 & 25/11 & \textcolor{RaciR}{\textbf{0.00}} \\
\rowcolor{RaciR!8}\textbf{5.1} & Manuales & 25/11 & 02/12 & 25/11 & 02/12 & \textcolor{RaciR}{\textbf{0.00}} \\
\rowcolor{RaciR!8}\textbf{5.2} & Informe final y exposición & 02/12 & 14/12 & 02/12 & 14/12 & \textcolor{RaciR}{\textbf{0.00}} \\
\rowcolor{RaciR!8}\textbf{5.3} & Lecciones aprendidas y acta de cierre & 14/12 & 16/12 & 14/12 & 16/12 & \textcolor{RaciR}{\textbf{0.00}} \\
\end{longtable}
}

\vspace{0.1cm}
\noindent\textbf{Actividades críticas (holgura cero):} 2.1, 1.2, 2.2, 2.3, 3.1, 3.2, 3.3, 3.4, 3.6, 4.3, 4.4, 5.1, 5.2, 5.3.

\vspace{0.1cm}
\noindent La ruta crítica dura \textbf{15.60 semanas} de calendario y el proyecto termina el \textbf{16/12/2026} (semana 17), antes del fin de la semana 17 (18/12/2026). La holgura del proyecto es de 0.40 semanas.

\begin{itemize}[leftmargin=*]
    \item Desviación estándar de la ruta crítica: $\sigma = \sqrt{\sum \sigma_i^2} = 0.60$ semanas.
    \item Probabilidad de terminar hasta el viernes 18/12/2026 (fin de la semana 17): $Z = (18.00 - 17.60)/0.60 = 0.67$, es decir \textbf{75\,\%}.
    \item Probabilidad de terminar hasta el jueves 24/12/2026 (cierre del semestre): $Z = 2.00$, es decir \textbf{98\,\%}.
\end{itemize}

\Needspace{20\baselineskip}
\section{Diagrama de Gantt por Semanas}
\begin{center}
\begin{tikzpicture}[font=\tiny]
\fill[AzulClaro] (0,0.05) rectangle (10.30,0.55);
\node[font=\tiny\bfseries, text=AzulUNA] at (0.30,0.30) {2};
\node[font=\tiny\bfseries, text=AzulUNA] at (0.91,0.30) {3};
\node[font=\tiny\bfseries, text=AzulUNA] at (1.51,0.30) {4};
\node[font=\tiny\bfseries, text=AzulUNA] at (2.12,0.30) {5};
\node[font=\tiny\bfseries, text=AzulUNA] at (2.73,0.30) {6};
\node[font=\tiny\bfseries, text=AzulUNA] at (3.33,0.30) {7};
\node[font=\tiny\bfseries, text=AzulUNA] at (3.94,0.30) {8};
\node[font=\tiny\bfseries, text=AzulUNA] at (4.54,0.30) {9};
\node[font=\tiny\bfseries, text=AzulUNA] at (5.15,0.30) {10};
\node[font=\tiny\bfseries, text=AzulUNA] at (5.76,0.30) {11};
\node[font=\tiny\bfseries, text=AzulUNA] at (6.36,0.30) {12};
\node[font=\tiny\bfseries, text=AzulUNA] at (6.97,0.30) {13};
\node[font=\tiny\bfseries, text=AzulUNA] at (7.57,0.30) {14};
\node[font=\tiny\bfseries, text=AzulUNA] at (8.18,0.30) {15};
\node[font=\tiny\bfseries, text=AzulUNA] at (8.79,0.30) {16};
\node[font=\tiny\bfseries, text=AzulUNA] at (9.39,0.30) {17};
\node[font=\tiny\bfseries, text=AzulUNA] at (10.00,0.30) {18};
\fill[GrisFila!60!gray!25] (4.24,0.05) rectangle (4.85,-11.04);
\fill[GrisFila!60!gray!25] (9.09,0.05) rectangle (9.69,-11.04);
\fill[RaciR!6] (9.69,0.05) rectangle (10.30,-11.04);
\draw[gray!25] (0.00,0.05) -- (0.00,-11.04);
\draw[gray!25] (0.61,0.05) -- (0.61,-11.04);
\draw[gray!25] (1.21,0.05) -- (1.21,-11.04);
\draw[gray!25] (1.82,0.05) -- (1.82,-11.04);
\draw[gray!25] (2.42,0.05) -- (2.42,-11.04);
\draw[gray!25] (3.03,0.05) -- (3.03,-11.04);
\draw[gray!25] (3.64,0.05) -- (3.64,-11.04);
\draw[gray!25] (4.24,0.05) -- (4.24,-11.04);
\draw[gray!25] (4.85,0.05) -- (4.85,-11.04);
\draw[gray!25] (5.45,0.05) -- (5.45,-11.04);
\draw[gray!25] (6.06,0.05) -- (6.06,-11.04);
\draw[gray!25] (6.66,0.05) -- (6.66,-11.04);
\draw[gray!25] (7.27,0.05) -- (7.27,-11.04);
\draw[gray!25] (7.88,0.05) -- (7.88,-11.04);
\draw[gray!25] (8.48,0.05) -- (8.48,-11.04);
\draw[gray!25] (9.09,0.05) -- (9.09,-11.04);
\draw[gray!25] (9.69,0.05) -- (9.69,-11.04);
\draw[gray!25] (10.30,0.05) -- (10.30,-11.04);
\node[anchor=east, font=\tiny\bfseries, text=Rama1] at (-0.1,-0.40) {1 Gestión del proyecto};
\node[anchor=east, font=\tiny, text width=4.7cm, align=right] at (-0.1,-0.80) {1.1 Gobernanza del equipo};
\fill[Rama1] (0.00,-0.92) rectangle (1.21,-0.68);
\draw[white, line width=0.4pt] (0.00,-0.80) -- (1.21,-0.80);
\node[anchor=east, font=\tiny, text width=4.7cm, align=right] at (-0.1,-1.20) {1.2 Viabilidad y Acta de Constitución};
\fill[RaciR] (0.61,-1.32) rectangle (1.21,-1.08);
\draw[white, line width=0.4pt] (0.61,-1.20) -- (1.21,-1.20);
\node[anchor=east, font=\tiny, text width=4.7cm, align=right] at (-0.1,-1.60) {1.3 Planificación base};
\fill[Rama1] (1.21,-1.72) rectangle (1.85,-1.48);
\node[anchor=east, font=\tiny, text width=4.7cm, align=right] at (-0.1,-2.00) {1.4 Planes de soporte};
\fill[Rama1] (1.85,-2.12) rectangle (2.48,-1.88);
\node[anchor=east, font=\tiny, text width=4.7cm, align=right] at (-0.1,-2.40) {1.5 Seguimiento y control};
\fill[pattern=north east lines, pattern color=GrisOscuro!60] (2.48,-2.51) rectangle (9.45,-2.29);
\draw[GrisOscuro!60] (2.48,-2.51) rectangle (9.45,-2.29);
\node[anchor=east, font=\tiny, text width=4.7cm, align=right] at (-0.1,-2.80) {1.6 Entregas de la Unidad II};
\fill[Rama1] (3.76,-2.92) rectangle (4.90,-2.68);
\node[anchor=east, font=\tiny\bfseries, text=Rama2] at (-0.1,-3.20) {2 Análisis y diseño};
\node[anchor=east, font=\tiny, text width=4.7cm, align=right] at (-0.1,-3.60) {2.1 Ficha técnica y requerimientos};
\fill[RaciR] (0.00,-3.72) rectangle (0.61,-3.48);
\draw[white, line width=0.4pt] (0.00,-3.60) -- (0.61,-3.60);
\node[anchor=east, font=\tiny, text width=4.7cm, align=right] at (-0.1,-4.00) {2.2 Modelo de datos y arquitectura};
\fill[RaciR] (1.21,-4.12) rectangle (1.85,-3.88);
\node[anchor=east, font=\tiny, text width=4.7cm, align=right] at (-0.1,-4.40) {2.3 Diseño de pantallas};
\fill[RaciR] (1.21,-4.52) rectangle (1.85,-4.28);
\node[anchor=east, font=\tiny\bfseries, text=Rama3] at (-0.1,-4.80) {3 Construcción del software};
\node[anchor=east, font=\tiny, text width=4.7cm, align=right] at (-0.1,-5.20) {3.1 Motor CASA SCORE};
\fill[RaciR] (1.21,-5.32) rectangle (2.48,-5.08);
\node[anchor=east, font=\tiny, text width=4.7cm, align=right] at (-0.1,-5.60) {3.2 Publicación de cuartos};
\fill[RaciR] (1.85,-5.72) rectangle (2.48,-5.48);
\node[anchor=east, font=\tiny, text width=4.7cm, align=right] at (-0.1,-6.00) {3.3 Ranking y detalle (MVP Unidad I)};
\fill[RaciR] (2.48,-6.12) rectangle (3.76,-5.88);
\node[anchor=east, font=\tiny, text width=4.7cm, align=right] at (-0.1,-6.40) {3.4 Nube, registro e inicio de sesión};
\fill[RaciR] (3.76,-6.52) rectangle (6.27,-6.28);
\node[anchor=east, font=\tiny, text width=4.7cm, align=right] at (-0.1,-6.80) {3.5 Punto de referencia y nuevos criterios};
\fill[Rama3] (3.76,-6.92) rectangle (5.63,-6.68);
\node[anchor=east, font=\tiny, text width=4.7cm, align=right] at (-0.1,-7.20) {3.6 Perfil y mis cuartos};
\fill[RaciR] (6.27,-7.32) rectangle (6.91,-7.08);
\node[anchor=east, font=\tiny\bfseries, text=Rama4] at (-0.1,-7.60) {4 Pruebas y validación con clientes};
\node[anchor=east, font=\tiny, text width=4.7cm, align=right] at (-0.1,-8.00) {4.1 Pruebas del motor};
\fill[Rama4] (2.48,-8.12) rectangle (3.12,-7.88);
\node[anchor=east, font=\tiny, text width=4.7cm, align=right] at (-0.1,-8.40) {4.2 Levantamiento de cuartos};
\fill[Rama4] (1.85,-8.52) rectangle (2.48,-8.28);
\node[anchor=east, font=\tiny, text width=4.7cm, align=right] at (-0.1,-8.80) {4.3 Piloto con clientes};
\fill[RaciR] (6.27,-8.92) rectangle (6.91,-8.68);
\node[anchor=east, font=\tiny, text width=4.7cm, align=right] at (-0.1,-9.20) {4.4 Pruebas de aceptación};
\fill[RaciR] (6.91,-9.32) rectangle (7.54,-9.08);
\node[anchor=east, font=\tiny\bfseries, text=Rama5] at (-0.1,-9.60) {5 Cierre};
\node[anchor=east, font=\tiny, text width=4.7cm, align=right] at (-0.1,-10.00) {5.1 Manuales};
\fill[RaciR] (7.54,-10.12) rectangle (8.18,-9.88);
\node[anchor=east, font=\tiny, text width=4.7cm, align=right] at (-0.1,-10.40) {5.2 Informe final y exposición};
\fill[RaciR] (8.18,-10.52) rectangle (9.13,-10.28);
\node[anchor=east, font=\tiny, text width=4.7cm, align=right] at (-0.1,-10.80) {5.3 Lecciones aprendidas y acta de cierre};
\fill[RaciR] (9.13,-10.92) rectangle (9.45,-10.68);
\draw[RaciR, dashed, thick] (9.69,0.6) -- (9.69,-11.04);
\node[anchor=south, text=RaciR, font=\tiny\bfseries] at (9.69,0.55) {fin de la semana 17};
\end{tikzpicture}
\\[2pt]
{\scriptsize \textcolor{RaciR}{\rule{9pt}{5pt}}~Ruta crítica \quad \textcolor{Rama1}{\rule{9pt}{5pt}}~Rama 1 \quad \textcolor{Rama2}{\rule{9pt}{5pt}}~Rama 2 \quad \textcolor{Rama3}{\rule{9pt}{5pt}}~Rama 3 \quad \textcolor{Rama4}{\rule{9pt}{5pt}}~Rama 4 \quad \textcolor{Rama5}{\rule{9pt}{5pt}}~Rama 5 \quad \tikz{\fill[pattern=north east lines, pattern color=GrisOscuro!60] (0,0) rectangle (0.32,0.18);}~Nivel de esfuerzo \quad \tikz{\fill[GrisFila!60!gray!25] (0,0) rectangle (0.32,0.18);}~Semana de exámenes}
\end{center}

\Needspace{12\baselineskip}
\section{Línea Base del Cronograma}
{\small
\rowcolors{2}{white}{GrisFila}
\noindent\begin{tabularx}{\textwidth}{@{}L C{0.9cm} C{2.75cm} C{2.75cm} C{2.3cm}@{}}
\toprule
\rowcolor{AzulClaro}
\textbf{Hito} & \textbf{Act.} & \textbf{Fecha del Acta} & \textbf{Línea base} & \textbf{Diferencia} \\
\midrule
Acta de Constitución aprobada & 1.2 & 11/09 (sem.\,3) & 11/09 (sem.\,3) & \textcolor{Rama4}{\textbf{a tiempo}} \\
MVP de la Unidad I & 3.3 & 16/10 (sem.\,8) & 12/10 (sem.\,8) & \textcolor{Rama4}{\textbf{4 días antes}} \\
Plan de Dirección del Proyecto & 1.6 & 26/10 (sem.\,10) & 26/10 (sem.\,10) & \textcolor{Rama4}{\textbf{a tiempo}} \\
Piloto con estudiantes y propietarios & 4.3 & 20/11 (sem.\,13) & 17/11 (sem.\,13) & \textcolor{Rama4}{\textbf{3 días antes}} \\
MVP final aceptado & 4.4 & 11/12 (sem.\,16) & 25/11 (sem.\,14) & \textcolor{Rama4}{\textbf{16 días antes}} \\
Informe final y exposición & 5.2 & 14/12 (sem.\,17) & 14/12 (sem.\,17) & \textcolor{Rama4}{\textbf{a tiempo}} \\
Acta de cierre & 5.3 & 18/12 (sem.\,17) & 16/12 (sem.\,17) & \textcolor{Rama4}{\textbf{2 días antes}} \\
\bottomrule
\end{tabularx}
}

\vspace{0.2cm}
\noindent La línea base del cronograma se aprueba con esta versión. Durante la ejecución:
\begin{itemize}[leftmargin=*]
    \item Cada martes se registra el avance real de cada actividad y se compara con esta línea base (variación del cronograma e índice SPI).
    \item Un retraso mayor a 3 días en una actividad de la ruta crítica activa una acción correctiva en la reunión de control.
    \item Un cambio de fecha de un hito requiere una solicitud de cambio aprobada por el equipo.
\end{itemize}

\section{Firmas}
\noindent En señal de conformidad con el presente cronograma, los integrantes del equipo firman a continuación:

\vspace{1.5cm}
\begin{center}
\begin{tabular}{c c c}
    \parbox{6.5cm}{\centering \rule{6.0cm}{0.6pt}\\[5pt] \textbf{Clyde Neil Paricahua Pari}\\[2pt] {\footnotesize Project Manager}}
    & \hspace{0.8cm} &
    \parbox{6.5cm}{\centering \rule{6.0cm}{0.6pt}\\[5pt] \textbf{Ruth Karina Apaza Solis}\\[2pt] {\footnotesize Analista de Planificación}}
    \\[2.0cm]
    \parbox{6.5cm}{\centering \rule{6.0cm}{0.6pt}\\[5pt] \textbf{Eliseo Tarqui Ajahuana}\\[2pt] {\footnotesize Especialista Técnico / Arquitectura}}
    & \hspace{0.8cm} &
    \parbox{6.5cm}{\centering \rule{6.0cm}{0.6pt}\\[5pt] \textbf{Viviana Mary Quisbert Quispe}\\[2pt] {\footnotesize Gestora de Comunicación, Riesgos y Documentación}}
\end{tabular}
\end{center}
\end{document}
